\documentclass[UTF8,zihao=-4,a4paper]{ctexart}

\usepackage{geometry}
\geometry{left=24mm,right=24mm,top=25mm,bottom=25mm,headheight=24pt}
\usepackage{xcolor}
\usepackage{fontspec}
\usepackage{setspace}
\usepackage{titlesec}
\usepackage{titling}
\usepackage{fancyhdr}
\usepackage{hyperref}
\usepackage{enumitem}
\usepackage[most]{tcolorbox}
\usepackage{tikz}
\usepackage{eso-pic}
\usepackage{needspace}
\usepackage{microtype}

\setmainfont{Avenir Next}
\setsansfont{Avenir Next}
\setCJKmainfont{Songti SC}
\setCJKsansfont{PingFang SC}
\setCJKmonofont{PingFang SC}

\definecolor{Ink}{HTML}{172033}
\definecolor{Muted}{HTML}{5F6C82}
\definecolor{Primary}{HTML}{2563EB}
\definecolor{PrimaryDark}{HTML}{1E3A8A}
\definecolor{Accent}{HTML}{7C3AED}
\definecolor{SoftBlue}{HTML}{EFF6FF}
\definecolor{SoftViolet}{HTML}{F5F3FF}
\definecolor{SoftAmber}{HTML}{FFFBEB}
\definecolor{Rule}{HTML}{DBEAFE}
\definecolor{Paper}{HTML}{FBFDFF}

\pagecolor{Paper}
\color{Ink}
\onehalfspacing
\setlength{\parindent}{2em}
\setlength{\parskip}{0.28em}

\hypersetup{
  colorlinks=true,
  linkcolor=PrimaryDark,
  urlcolor=Primary,
  pdftitle={AI 时代，我发现自己快不会思考了},
  pdfauthor={李泉胜},
  pdfsubject={Personal Blog},
  pdfkeywords={AI, 思考, 科研, 写作, Vibe Coding}
}

\pagestyle{fancy}
\fancyhf{}
\fancyhead[L]{\small\sffamily\color{Muted}李泉胜 · Personal Blog}
\fancyhead[R]{\small\sffamily\color{Muted}AI 时代，我发现自己快不会思考了}
\fancyfoot[C]{\small\sffamily\color{Muted}\thepage}
\renewcommand{\headrulewidth}{0.4pt}
\renewcommand{\headrule}{\hbox to\headwidth{\color{Rule}\leaders\hrule height \headrulewidth\hfill}}
\renewcommand{\footrulewidth}{0pt}

\titleformat{\section}
  {\Needspace{6\baselineskip}\Large\bfseries\sffamily\color{PrimaryDark}}
  {\thesection}{0.7em}{}
  [\vspace{0.2em}{\color{Primary}\titlerule[1.1pt]}]
\titleformat{\subsection}
  {\large\bfseries\sffamily\color{Ink}}
  {\thesubsection}{0.6em}{}
\titlespacing*{\section}{0pt}{1.8em}{0.9em}
\titlespacing*{\subsection}{0pt}{1.1em}{0.4em}

\setlist[itemize]{leftmargin=2em,itemsep=0.2em,topsep=0.3em}
\setlist[enumerate]{leftmargin=2em,itemsep=0.2em,topsep=0.3em}

\newtcolorbox{leadbox}{
  enhanced,
  breakable,
  colback=SoftBlue,
  colframe=Primary!28,
  boxrule=0.8pt,
  arc=3mm,
  left=4mm,
  right=4mm,
  top=3mm,
  bottom=3mm,
  borderline west={3pt}{0pt}{Primary},
}

\newtcolorbox{notebox}[1]{
  enhanced,
  breakable,
  title={#1},
  colback=SoftViolet,
  colframe=Accent!30,
  coltitle=PrimaryDark,
  fonttitle=\bfseries\sffamily,
  boxrule=0.7pt,
  arc=3mm,
  left=4mm,
  right=4mm,
  top=3mm,
  bottom=3mm,
}

\newtcolorbox{amberbox}[1]{
  enhanced,
  breakable,
  title={#1},
  colback=SoftAmber,
  colframe=orange!28,
  coltitle=Ink,
  fonttitle=\bfseries\sffamily,
  boxrule=0.7pt,
  arc=3mm,
  left=4mm,
  right=4mm,
  top=3mm,
  bottom=3mm,
}

\newcommand{\keyword}[1]{\textcolor{PrimaryDark}{\bfseries #1}}
\newcommand{\divider}{\vspace{0.6em}{\color{Rule}\hrule}\vspace{0.6em}}

\begin{document}
\pagenumbering{gobble}

\begin{titlepage}
\thispagestyle{empty}
\AddToShipoutPictureBG*{%
\begin{tikzpicture}[remember picture,overlay]
  \fill[SoftBlue] (current page.north west) rectangle ([xshift=17mm]current page.south west);
  \fill[Primary] ([xshift=11mm]current page.north west) rectangle ([xshift=17mm]current page.south west);
  \fill[Accent!18] ([xshift=-48mm,yshift=-18mm]current page.north east) circle (36mm);
  \fill[Primary!12] ([xshift=20mm,yshift=18mm]current page.south west) circle (42mm);
\end{tikzpicture}}
\vspace*{20mm}
\begin{flushleft}
{\sffamily\bfseries\color{Primary}\Large PERSONAL BLOG · 001}\par
\vspace{18mm}
{\sffamily\bfseries\fontsize{30}{39}\selectfont AI 时代，\\我发现自己快不会思考了}\par
\vspace{10mm}
{\color{Muted}\Large 一次关于工具依赖、科研训练与自我夺回的复盘}\par
\vspace{18mm}
\begin{leadbox}
\large
我不是要少用 AI。相反，我以后一定还会更深地使用 AI。只是我越来越清楚：真正重要的不是让 AI 替我完成任务，而是让我重新回到问题定义、结构搭建、原理理解和结果审查之中。
\end{leadbox}
\vfill
{\sffamily\large 李泉胜}\par
\vspace{2mm}
{\sffamily\color{Muted}中国科学院上海技术物理研究所 · Personal Academic Homepage}\par
\vspace{2mm}
{\sffamily\color{Muted}发表日期：2026年5月11日}\par
\end{flushleft}
\end{titlepage}

\clearpage
\pagenumbering{roman}
\tableofcontents
\clearpage
\pagenumbering{arabic}

\section*{写在前面}
\addcontentsline{toc}{section}{写在前面}
最近我越来越强烈地感觉到一件事：我好像越来越会用 AI 了，但也越来越不会自己思考了。

这句话听起来有点矛盾。因为从表面上看，AI 确实让我变“厉害”了。

我能更快写论文、写申报书、做项目方案、改代码、整理资料、生成图表思路，甚至可以快速搭出一个看起来很完整的科研工具原型。外人看起来可能会觉得：这个人好像什么都会一点，写作、科研、代码、产品、自动化工具都能碰。

但我自己很清楚，很多东西只是“碰过”，不是真的变成了能力。

有些时候，我用 AI 做完一份申报书，或者做完一篇论文初稿，回头再看，甚至会有一种很空的感觉：

\begin{leadbox}
这东西到底是不是我想出来的？\par
里面的逻辑我真的掌握了吗？\par
如果别人追问其中一个结构为什么这么安排，我能不能讲清楚？
\end{leadbox}

很多时候，答案并不乐观。

我现在越来越觉得，AI 最大的风险，不是替代执行，而是诱导我把“本该自己完成的判断”也外包出去。

\section{我的现状：信息太多，工具太多，判断却越来越少}
我现在最大的问题，不是没有资源，也不是没有工具，而是工具和信息太多了。

每天都有新的模型、新的 AI 编程工具、新的科研自动化工具、新的提示词、新的工作流。看起来每一个都值得学，每一个都像是下一代生产力入口。

我什么都想学。

新模型出来了，我想试；新产品上线了，我想试；别人说某个工具很强，我也想马上去试。表面上看，我好像一直站在 AI 信息的前沿，但实际上，我并没有因此形成稳定的产出。

更具体一点说，我很容易陷入一种“工具猎奇”和“资源搜集”的兴奋感里。

哪里能白嫖模型，哪里有便宜订阅，哪里有中转站，哪里有新的反代方案，哪里又出现了新的替代入口，我都想去看看。我会去 Telegram 群、NodeSeek、Linux.do，还有各种 AI 社群和信息流里刷最新消息。

一开始，我会觉得这是一种信息优势。

我知道得比别人早。\par
我试过的工具比别人多。\par
我能找到一些别人不知道的入口。\par
我还能把这些信息分享给别人，好像自己很懂 AI，很有资源，很强。

但冷静下来想，这些东西真的重要吗？

除非我真的要把它变成一个长期项目，否则这些入口和渠道本身并不重要。它们最多只能让我短时间内用上某个工具，或者省下一点钱。

我知道得再早，也不代表我能产出。\par
我试得再多，也不代表我有判断。\par
我收藏再多入口，也不代表我形成了自己的长期路线。

更糟糕的是，我还会为了这些信息熬夜。

有时候晚上本来应该休息，结果又去看 Telegram 群里有没有新消息，NodeSeek 上有没有新讨论，Linux.do 上有没有新的教程、工具或者账号信息。看起来像是在学习，像是在追前沿，但很多时候只是被信息流牵着走。

刷到最后，大脑是乱的。\par
注意力是散的。\par
睡眠也被打乱了。\par
第二天真正要写论文、看文献、做项目的时候，反而没有精力了。

还有一个很明显的问题是，我越来越习惯于“复制粘贴式”的 AI 使用方式。

看到一段材料，复制进去；\par
想要一段文字，让 AI 生成；\par
需要一个标题，让 AI 生成；\par
需要一张图，也让 AI 生成。

很多时候就是 Ctrl+C、Ctrl+V，再加一句提示词。内容很快出来了，图片也很快生成了，任务好像也在不断推进。

但这种方式用多了之后，人会越来越习惯跳过自己的加工过程。

AI 帮我生成了一段话，不代表我真的懂这段话。\par
AI 帮我生成了一张图，不代表我真的理解这张图背后的逻辑。\par
AI 帮我整理出一个结构，不代表我真的知道为什么要这样安排。

\begin{notebox}{三种失控}
\begin{itemize}
  \item \keyword{输入失控}：信息太多，工具太多，注意力被不断切碎。
  \item \keyword{过程失控}：越来越习惯复制粘贴、生成内容、让 AI 代替我往前推进。
  \item \keyword{判断失控}：结果看起来完成了，但我不知道它为什么这样完成，也不知道它到底是不是我真正想要的东西。
\end{itemize}
\end{notebox}

所以我不是不用 AI。我也不是否定工具。

我只是越来越清楚地发现，如果我只是在追工具、刷信息、找入口、试模型，却没有稳定产出，也没有真正参与核心判断，那我并没有变强。

我只是变得更忙、更散、更依赖 AI 了。

\section{我为什么想改变：我到了一个必须重新训练自己的阶段}
我现在正好处在一个转折点。

本科即将结束，马上进入研究生阶段。我原本以为，自己可以靠 AI 大展身手。毕竟我已经比很多人更早接触这些工具，也花了很多时间研究 AI 写作、AI 编程、AI 科研工作流。

但现实给我的感觉是：没有那么容易大展身手。

因为会用 AI 这件事，正在迅速变得不稀缺。

今天我会用 Codex，别人明天也会用。\par
我会用 Claude Code，别人也会用。\par
我会用 AI 写论文、写代码、做方案，等这些工具普及以后，所有人都会用。

那我的壁垒在哪里？

如果我的能力只是“比别人更早知道一个工具”，那这个壁垒是很脆弱的。因为工具一旦普及，差距很快就会被抹平。

我真正担心的是：

我以为自己在用 AI 提升能力，实际上只是把任务外包给 AI；\par
我以为自己在做科研，实际上只是让 AI 拼接科研语言；\par
我以为自己在做项目，实际上只是让 AI 替我搭一个看起来很完整的壳。

这对我来说很危险。

因为我真正想要的，不是简单地“会用 AI”，而是用 AI 放大我的专业能力。

我的方向是红外探测器。这个领域本身就不是靠几段漂亮的学术语言能撑起来的。器件结构、物理机制、读出电路、测试方法、性能参数、文献脉络，每一层都需要真正理解。

如果我自己不参与论文结构的思考，不亲自判断一个创新点是否合理，不去推敲一段论证为什么成立，只是让 AI 自动代替我完成，那么最后产出的东西大概率是空的。

它可能语言很顺，结构很完整，格式也像论文。但它不一定真的有科研判断。

而科研里最核心的东西，恰恰就是判断。

这种担心不是凭空来的，而是我已经在本科毕业论文里真实经历过一次。

坦白说，那段时间我对 AI 的依赖已经非常严重了。

不只是润色，不只是改语句，也不只是查资料。很多正文、结构、图示思路、Word 排版，甚至答辩 PPT，都是 AI 帮我一步步完成的。

当时我当然觉得很爽。

论文推进得很快。\par
PPT 做得也快。\par
很多原本让我头疼的内容，被 AI 一下子整理得很完整。

但现在回头看，我反而有点后怕。

因为我发现，文件是完成了，但我对整个过程的掌控感并不强。很多内容看起来是我的成果，但里面有多少是真正经过我自己思考、判断、组织和验证的？这个问题我没法轻轻松松回答。

毕业论文不应该只是一个最终文件。它本来应该是本科阶段专业训练的一次总结，是我把知识、设计、分析、图表、表达和答辩逻辑重新组织起来的过程。

可如果这个过程里，大量关键环节都由 AI 替我完成，那我到底训练到了什么？

我不想到了研究生阶段，还继续用这种方式做科研。\par
不想论文写完了，却不知道自己到底写了什么。\par
不想 PPT 做完了，却只是照着 AI 给的结构去讲。\par
更不想把本来应该训练自己的机会，全部交给 AI。

\begin{amberbox}{一句提醒}
AI 可以帮我，但不能替我完成成长。
\end{amberbox}

\section{大部分人会怎么用 AI？我觉得问题在哪里？}
我观察到很多人使用 AI，其实有几种常见方式。这些方式，我自己也都经历过。

所以与其说我是在批评别人，不如说我是在复盘以前的自己。

第一种，是把 AI 当成聊天搜索框。

问一句，答一句。让它解释概念、翻译文字、润色段落、写个提纲。

这种方式不算错，但效率并不高。它只是把 AI 当成一个更聪明的搜索引擎。

不过我觉得这不是最大的问题。因为随着 AI 工具越来越普及，工作流工具越来越成熟，只停留在对话端的问题迟早会被解决。

更大的问题是第二种：把 AI 当成全流程代工。

给 AI 一个模糊目标，然后让它从头做到尾。自己不监督，不拆解，不检查，不判断。最后拿到一个看似完整的结果，就以为任务完成了。

不错，这就是我。

这种“全流程代工”的问题，在我身上主要体现在两个场景里。

\subsection{科研：完成不等于掌握}
我会让 AI 直接写一篇文章，直接规划一个项目，直接生成一份申报书，直接搭一个论文框架。它写得越完整，我越容易产生一种错觉：好像我真的完成了这件事。

尤其是做论文、综述、申报材料的时候，AI 很容易给出一套看起来很顺的结构。标题完整，逻辑完整，语言也很像那么回事。它甚至能帮我把背景意义、研究内容、技术路线、创新点、预期成果全部写出来。

但后来我发现，完成不等于掌握。

AI 给出的结果越完整，人反而越容易跳过中间的思考。而真正的科研能力，往往就在中间过程里。

比如：

为什么这样分章节？\par
为什么这个概念要放在前面？\par
为什么这个实验能支撑这个结论？\par
为什么这个创新点不是表面创新？\par
为什么这个方法在我的领域里有意义？

这些问题如果我不问，AI 不会主动替我形成真正的能力。

\subsection{编程：消费“做出东西”的快感}
第二个场景，是编程。

今年很火的 Vibe Coding，其实也有类似的问题。

我并不反对 Vibe Coding。相反，我觉得它确实降低了做东西的门槛。很多原本只停留在脑子里的想法，现在可以很快变成一个能看的 Demo。少部分项目如果继续打磨，也确实可能变成有意义的东西。

但问题是，有些时候，人会为了 Vibe Coding 而 Vibe Coding。

这也是我。

突然想到一个小项目，就马上让 AI 帮我写代码、搭页面、修 bug、推送到 GitHub。整个过程很快，几小时就能做出一个看起来还不错的东西。页面能跑，功能能演示，README 也像模像样，甚至部署链接也能发出去。

然后我会很自然地把它发到小红书，安利自己的 AI Coding 成果。

如果有人点赞、收藏，甚至 GitHub 上有几个 Star，我会很开心。那一瞬间确实会觉得：我好像很厉害，我好像真的做出了一个产品。

但冷静下来想，很多时候我只是做出了一个“看起来像产品”的东西。

一问技术细节，我可能讲不清楚。\par
代码为什么这样组织，我不一定知道。\par
这个功能为什么这样实现，我不一定理解。\par
GitHub 推送、部署、报错修复，很多也是 AI 帮我一步步完成的。

表面上看，我是在做项目。实际上，有时候我只是在消费“做出东西”的快感。

这和追新模型、刷 AI 工具、找白嫖入口，本质上有点像：它们都能让我短时间内产生一种“我在进步”的感觉。但如果没有复盘、没有理解、没有持续迭代，它们很难真正沉淀成能力。

这两个场景看起来不一样，一个是科研，一个是编程，但底层问题其实是一样的：

\begin{leadbox}
我把过程交给了 AI。\par
我跳过了真正需要自己思考的部分。\par
我拿到了一个看似完整的结果，却没有真正形成能力。
\end{leadbox}

难点在于，这些方式真的很爽。

AI 代写很快，Vibe Coding 很快，发出去也很快。人能马上看到结果，马上获得反馈，马上觉得自己在进步。

但真正的思考、复盘、验证和长期积累，反馈都很慢。

所以问题不是大家不知道深度思考重要，而是即时反馈太诱人，长期训练太慢。

这才是 AI 时代最危险的地方：不是不会用 AI，而是误以为“让 AI 完成任务”就等于“自己具备能力”。

\section{我想提出的方案：不是让 AI 主导我，而是我主导 AI}
我现在越来越明确一个观点：

\begin{leadbox}
AI 不应该主导我。\par
我应该主导 AI。
\end{leadbox}

我不想再把 AI 当成一个替我思考的机器，而是要把它当成一个被我调度、被我审查、被我约束的外部能力。

也就是说，我应该成为 AI 的审查员，而不是 AI 的跟随者。

在我看来，AI 时代真正重要的能力，不是单纯会不会提问，也不是会不会写提示词，而是能不能搭建一套属于自己的思考链。

\begin{notebox}{四步思考链}
\begin{enumerate}
  \item \keyword{问题必须由我来定义}：我得知道我到底想解决什么问题，而不是把一个模糊目标扔给 AI，让它替我决定方向。
  \item \keyword{核心结构必须由我来搭建}：论文的章节结构、项目的技术路线、研究的逻辑链条，这些东西不能完全交给 AI。
  \item \keyword{底层原理必须由我自己理解}：尤其在科研里，如果我自己不懂器件原理、测试逻辑和参数关系，漂亮的话就只是空壳。
  \item \keyword{结果必须由我来审查}：我要判断它有没有偷换概念，有没有逻辑断裂，有没有为了完整而完整。
\end{enumerate}
\end{notebox}

简单说，就是：

针对输入失控，我要减少无目的追新，只保留能服务长期产出的工具。\par
针对过程失控，我要把 AI 放在辅助位置，而不是让它从头到尾代工。\par
针对判断失控，我要在每次使用 AI 后追问：它为什么这样写？为什么这样做？哪里可能是错的？

我现在更愿意用一个比喻来理解这件事：

地基应该由我来盖。AI 可以帮我盖地基外面的建筑。

它可以帮我把房子修得更快、更漂亮、更完整。但如果地基不是我打的，这个房子看起来再好，也不稳。

\section{我的尝试：先从几个小动作开始}
我现在还没有一个很成熟的方法论，但我想先从几个小动作开始。

\begin{amberbox}{四个小动作}
\begin{enumerate}
  \item \keyword{用 AI 之前，先自己写问题框架。} 哪怕只写几行，也要先把“我要解决什么问题”“我希望得到什么结果”“我自己的初步判断是什么”写下来。
  \item \keyword{让 AI 审查结构，而不是直接代写结构。} 写论文、写申报书、做项目方案时，先自己搭一版框架，再让 AI 帮我找漏洞、补缺口、提反例。
  \item \keyword{对 AI 生成的内容追问“为什么成立”。} 不只看语言顺不顺，还要问依据是什么、逻辑链能不能支撑、有没有偷换概念。
  \item \keyword{每次使用 AI 后做一个小复盘。} 追问哪些是 AI 完成的，哪些是我自己判断的，哪些地方只是接受了答案但没有理解。
\end{enumerate}
\end{amberbox}

这几个动作看起来很小，但对我来说很重要。

因为我不是要少用 AI。我以后一定还会更深地使用 AI。

但我必须把自己重新放回思考链条里。

\section{其他人的尝试}
这一部分我暂时不急着下结论。

后面我也想继续观察一些人的做法：有人把 AI 用成知识管理系统，有人用 AI 辅助科研阅读，有人用 AI 做产品原型，也有人只是不断追模型、追工具、追新入口。

这些做法里，有些可能真的值得长期参考，有些可能只是新的信息焦虑。

等我真正看到一些值得借鉴的案例，再单独写一篇。

\section*{结尾：我想重新夺回自己的思考}
\addcontentsline{toc}{section}{结尾：我想重新夺回自己的思考}
AI 时代最讽刺的一点是：

它明明可以帮助我们变得更强，但如果用错了，也可能让我们越来越懒得思考。

我现在并不是要减少使用 AI。

相反，我以后一定还会更深地使用 AI。

但我希望自己换一种方式使用它。

不是把所有东西都交给 AI。\par
不是让 AI 带着我走。\par
不是只追求更快完成任务。

而是让 AI 放大我的专业能力，放大我的判断，放大我的结构化思考。

我想真正学会的是：在 AI 很强的时代，怎样让自己也变强。

不是成为一个会调用工具的人。而是成为一个能定义问题、搭建结构、理解原理、审查结果的人。

最后我也想坦白一点：其实这篇文章本身，也是借助 AI 写出来的。

但这一次，问题是我自己提出来的。\par
痛点是我自己一点点梳理出来的。\par
文章结构是我先想清楚的。\par
哪些地方要真实，哪些地方不能写虚，哪些地方必须保留我的困惑，也是我自己判断的。

AI 只是帮我把这些零散的想法组织成更清楚的文字。

所以这篇文章本身，也算是一次练习：

不是让 AI 替我思考，而是我先思考，再让 AI 帮我表达。\par
不是让 AI 主导我，而是我来掌控 AI。

也许这才是我现在最想重新建立的关系。

\vspace{1em}
\begin{leadbox}
\centering\sffamily\bfseries\color{PrimaryDark}
在 AI 很强的时代，重新训练自己的判断力。
\end{leadbox}

\end{document}
